\section{Introduction}
\label{sec:intro}

\IEEEPARstart{F}{ederated} learning lets several institutions train one predictive model without exchanging records~\cite{mcmahan2017,kairouz2021}. In the cross-silo setting that matters for healthcare, the participants are a handful of sites, each behind its own firewall and each connected to the aggregation server by a link that the institution's network team controls~\cite{rieke2020,dayan2021}. The learning job is then one tenant among many on an uplink with a fixed capacity, a scheduled transfer window and an audit trail. Whether that link is an institutional WAN, a 5G/6G edge slice serving a mobile clinic, or a satellite backhaul serving a remote site, the operator asks how many bytes the job will push per window and whether it will fit~\cite{lim2020,niknam2020,chen2021joint}.

Prior work reduces communication either by shrinking each message (quantisation, sparsification, sketching~\cite{sattler2020,rothchild2020}) or by sending messages less often (local steps, adaptive periods, lazy uploads~\cite{stich2019,wang2019adacomm,chen2018lag}), and a separate line keeps part of the model local for personalisation~\cite{arivazhagan2019,collins2021,liang2020,oh2022,pillutla2022}. Layer-wise schedules that synchronise different parameter blocks at different frequencies~\cite{lee2023fedlama} sit between these lines. What is less studied is that the number of transmitting rounds is not only a network quantity. Under client-level differential privacy (DP)~\cite{mcmahan2018dp}, every transmitting round is a release that the privacy accountant must charge for and into which noise is injected. A communication quota therefore also sets a privacy cost, and the two can be traded jointly.

This article makes that coupling explicit and measures it. We cast the federated job as a tenant with a quota $\rho$ on the mean fraction of the shared model transmitted per round, and study how accuracy, network cost and DP budget respond when rounds are spent sparingly. We propose \tierfed{}, a scheduler that measures a per-tier dispersion proxy online and solves, exactly, the integer programme that minimises the drift term of a blockwise local-SGD bound subject to $\rho$, while keeping a per-site classifier head private. We compare it with full-rate personalisation (FedPer), drift correction (FedProx, SCAFFOLD), layer-wise adaptive intervals (a FedLAMA-style reimplementation), adaptive-period local SGD (AdaComm), lazy uploads (LAG) and fixed multi-rate schedules, all on one shared architecture table, with ten seeds, Holm-corrected paired tests and one-sided non-inferiority on both AUROC and AUPRC.

\textbf{Findings.} On a patient-deduplicated readmission benchmark:
\begin{itemize}
\item \emph{Synchronising less does not cost accuracy.} Every sparse schedule we evaluated, including fixed periods and \tierfed{}, matches or slightly exceeds full-rate FedPer; \tierfed{} at $\rho=0.25$ is non-inferior to FedPer on AUROC and AUPRC with \num{diabetesmlpTierfedUploadSavingPct}\% less upload (Section~\ref{sec:results-fixed}). An earlier version of this work reported the opposite for fixed schedules; that result came from patient leakage across splits and a scheduling bug, and we retract it.
\item \emph{Timing and size reductions compose.} With sparse ternary compression and error feedback, \tierfed{} uploads \num{compTierfedStcUploadMb}~MB per site instead of FedPer's \num{compFedperNoneUploadMb}~MB without loss of AUROC (Section~\ref{sec:results-comp}).
\item \emph{Fewer rounds buy privacy.} Under an identical per-client clipping DP mechanism and an exact R\'enyi accountant, round-sparse methods are charged a smaller $\varepsilon$ and retain more utility at the same noise multiplier; \tierfed{} dominates the full-rate methods at every noise level we tested and is level with AdaComm except at the extremes (Section~\ref{sec:results-privacy}).
\end{itemize}
We are equally explicit about the limits: a well-chosen fixed period and AdaComm reach comparable or better accuracy at comparable upload without DP; \tierfed{}'s divergence-driven differentiation between tiers is rarely active on this data; and neither update inversion nor membership inference separates the methods. The contribution is therefore an operator-facing scheduler with an exact budget solver, together with a careful, reproducible measurement of how the communication quota simultaneously sets network cost and privacy cost, rather than an accuracy improvement.

\textbf{Contributions.}
\begin{itemize}
\item A network-side formulation of cross-silo FL as a tenant with a communication quota, and a blockwise local-SGD bound whose drift term $\sum_g K_g^2\tau_{\max}^2(\Gamma_g+\sigma_g^2)$ motivates per-tier intervals (Sections~\ref{sec:system}, \ref{sec:analysis}).
\item \tierfed{}: an exact, sub-millisecond solver for the budget-constrained interval programme, driven by online dispersion measurements, composable with any compressor (Section~\ref{sec:method}).
\item A measurement of the communication--privacy coupling under a corrected client-level DP mechanism, against round-sparse and full-rate baselines (Section~\ref{sec:results-privacy}).
\item Network-level costing on five link profiles and an attack-surface evaluation on the slices each method actually transmits (Sections~\ref{sec:results-net}, \ref{sec:results-privacy}).
\item A reporting discipline: bit-identical reruns, patient-grouped splits, a single architecture table, Holm-corrected tests, non-inferiority on AUROC and AUPRC, every number generated from released logs, and an explicit list of configurations in which the method does not help.
\end{itemize}
